% !TeX program = pdflatex
\documentclass[aspectratio=169,utf8,t]{beamer}

% Idioma e Codificacao
\usepackage[brazilian]{babel}
\usepackage{booktabs}
\usepackage{array}
\usepackage{listings}
\usepackage{amsmath}
\usepackage{amssymb}
\usepackage{graphicx}
\usepackage{tikz}
\usetikzlibrary{shapes,arrows,positioning,calc}

% Carrega o Tema Modelo dos Scouts (Widescreen 16:9, Nunito Sans)
\usetheme{scouts}

% Definicao de linguagem JavaScript/JSON/MongoSH para o listings
\lstdefinelanguage{mongosh}{
  keywords={use, show, db, dbs, collections, insertOne, insertMany, find, findOne, updateOne, updateMany, deleteOne, deleteMany, bulkWrite, createCollection, runCommand, collMod, ObjectId, ISODate, NumberDecimal, Decimal128, true, false, null},
  sensitive=true,
  comment=[l]{//},
  morecomment=[s]{/*}{*/},
  morestring=[b]",
  morestring=[b]'
}

% Configuracao para Codigo-Fonte (MongoSH / Python / JSON)
\lstset{
    basicstyle=\fontsize{5.0pt}{6.1pt}\selectfont\ttfamily,
    keywordstyle=\bfseries\color{ScoutNavy},
    stringstyle=\color{ScoutGreen},
    commentstyle=\itshape\color{gray},
    backgroundcolor=\color{black!3},
    frame=single,
    rulecolor=\color{black!15},
    breaklines=true,
    breakatwhitespace=true,
    showstringspaces=false,
    tabsize=2,
    aboveskip=2pt,
    belowskip=1pt,
    literate={á}{{\'a}}1 {ã}{{\~a}}1 {â}{{\^a}}1 {à}{{\`a}}1
             {é}{{\'e}}1 {ê}{{\^e}}1 {í}{{\'i}}1 {ó}{{\'o}}1
             {õ}{{\~o}}1 {ô}{{\^o}}1 {ú}{{\'u}}1 {ç}{{\c{c}}}1
             {Á}{{\'A}}1 {Ã}{{\~A}}1 {É}{{\'E}}1 {Í}{{\'I}}1
             {Ó}{{\'O}}1 {Õ}{{\~O}}1 {Ú}{{\'U}}1 {Ç}{{\c{C}}}1
}

% Metadados Institucionais
\author{Prof. Me. Reginaldo Fernandes}
\institute{Instituto Federal do Ceará -- Campus Tauá \\ Tecnologia em Análise e Desenvolvimento de Sistemas}
\title{Bancos de Dados Não-Relacionais}
\subtitle{Aula 06: Do CRUD Básico e Python (\texttt{PyMongo}) à Modelagem de Documentos BSON}
\date{27 de Outubro de 2026}

\begin{document}

% ------------------------------------------------------------------------------
% Capa da Apresentacao
% ------------------------------------------------------------------------------
\changecolours[inverse]{ScoutGreen}
\frame[plain]{\titlepage}

% ------------------------------------------------------------------------------
% Sumario Executivo da Aula
% ------------------------------------------------------------------------------
\changecolours{ScoutGreen}
\begin{frame}{Agenda do Encontro}
  \vspace{0.3cm}
  \tableofcontents
\end{frame}

% ------------------------------------------------------------------------------
% Modulos da Aula
% ------------------------------------------------------------------------------
% !TeX root = main.tex
\changecolours{ScoutGreen}
\section[Fundamentos: RDBMS, Documentos e BSON]{Fundamentos Teóricos: Do RDBMS ao Modelo de Documentos e BSON}

% ------------------------------------------------------------------------------
\twocolframe[title={Do RDBMS ao Modelo Orientado a Documentos},leftcol=7.0cm,rightcol=7.0cm]{%
  \alert{Por que Normalizamos no SQL (1NF a 5NF)?}\par
  \itemseps[topsepi=0.2mm,itemsepi=0.5mm,labelsep=1.5mm,leftmargini=2.5mm]
  \scriptsize
  \begin{itemize}
    \item Em bancos relacionais, se guardarmos dados do aluno e matrículas na mesma tabela, surgem \textbf{redundâncias} e \textbf{anomalias de atualização} (\textit{Update Anomalies}): alterar o endereço do aluno exige atualizar várias linhas.
    \item As \textbf{Formas Normais (1NF--5NF)} resolvem isso fragmentando os dados em dezenas de tabelas ligadas por chaves estrangeiras.
    \item \textbf{O Custo:} Para montar uma única tela da aplicação, o banco precisa executar múltiplos \texttt{JOINs} pesados em disco!
  \end{itemize}
}{%
  \alert{A Filosofia do Modelo de Documentos}\par
  \itemseps[topsepi=0.2mm,itemsepi=0.5mm,labelsep=1.5mm,leftmargini=2.5mm]
  \scriptsize
  \begin{itemize}
    \item No MongoDB, a unidade básica não é uma linha plana, mas um \textbf{Documento BSON} rico que suporta até \textbf{100 níveis de aninhamento} e \textbf{16 MB}.
    \item \textbf{Modelagem Guiada pela Aplicação:} Em vez de fragmentar uma entidade em 5 tabelas, agrupamos os dados que pertencem ao mesmo objeto em subdocumentos e arrays.
    \item \textbf{Ganho Direto:} Leitura de objetos complexos em \textbf{1 única operação de I/O} e atomicidade nativa no documento!
  \end{itemize}
}

% ------------------------------------------------------------------------------
\begin{frame}{Terminologia SQL vs. MongoDB e Ecossistema de Ferramentas}
  \vspace{0.05cm}
  \makebox[\textwidth][c]{%
    \scriptsize
    \renewcommand{\arraystretch}{1.22}
    \begin{tabular}{p{2.2cm} p{3.4cm} p{6.4cm}}
      \toprule
      \textbf{Conceito SQL} & \textbf{MongoDB} & \textbf{Definição Técnica} \\
      \midrule
      \texttt{Database} & \texttt{Database} & Agrupador lógico (\texttt{mongodb\_books}, \texttt{ecommerce\_shopgrids}). \\
      \texttt{Table} & \texttt{Collection} & Conjunto de documentos com esquema dinâmico/flexível. \\
      \texttt{Row / Tuple} & \texttt{Document} (BSON) & Objeto serializado (até \textbf{16 MB} e \textbf{100 níveis}). \\
      \texttt{Column} & \texttt{Field} (Campo) & Par \texttt{chave: valor} (escalares, objetos e arrays). \\
      \texttt{Primary Key} & Campo \texttt{\_id} & Chave obrigatória e única (\texttt{ObjectId} de 12 bytes). \\
      \texttt{JOIN} & \textbf{Embedding} / \texttt{\$lookup} & Aninhamento direto ou agregação entre coleções. \\
      \bottomrule
    \end{tabular}%
  }
  \vspace{0.08cm}
  \centering{\scriptsize \textbf{Ferramentas Oficiais:} \textbf{1. MongoDB Compass} (GUI visual) $\bullet$ \textbf{2. \texttt{mongosh}} (Shell JS/Node.js) $\bullet$ \textbf{3. Drivers} (\texttt{PyMongo}).}
\end{frame}

% ------------------------------------------------------------------------------
\twocolframe[title={JSON vs. BSON (\textit{Binary JSON}): Os 3 Pilares de Projeto},leftcol=7.0cm,rightcol=7.0cm]{%
  \alert{Por que o MongoDB Usa BSON no Motor?}\par
  \itemseps[topsepi=0.2mm,itemsepi=0.5mm,labelsep=1.5mm,leftmargini=2.5mm]
  \scriptsize
  \begin{itemize}
    \item O JSON puro é texto: exige \textit{parsing} lento caractere por caractere e só possui um tipo numérico (\texttt{number} IEEE 754 float, onde $0.1 + 0.2 = 0.30000000000000004$).
    \item O formato \textbf{BSON} foi projetado sobre \textbf{3 características fundamentais}:
      \begin{enumerate}\scriptsize
        \item \textbf{\textit{Lightweight} (Leve):} Baixa sobrecarga espacial para tráfego de rede.
        \item \textbf{\textit{Traversable} (Percorrível):} Cabeçalhos com tamanho em bytes permitem saltar campos em $\mathcal{O}(1)$!
        \item \textbf{\textit{Efficient} (Eficiente):} Codificação e decodificação rápidas em tipos nativos de máquina (C/C++).
      \end{enumerate}
  \end{itemize}
}{%
  \alert{Principais Tipos de Dados BSON}\par
  \itemseps[topsepi=0.2mm,itemsepi=0.5mm,labelsep=1.5mm,leftmargini=2.5mm]
  \scriptsize
  \begin{itemize}
    \item \textbf{\texttt{ObjectId}:} Identificador binário de 12 bytes rápido de gerar e ordenado no tempo.
    \item \textbf{\texttt{String}:} Cadeia de caracteres UTF-8.
    \item \textbf{\texttt{Object} e \texttt{Array}:} Subdocumentos embutidos e listas de valores.
    \item \textbf{\texttt{Int32} e \texttt{Int64 (Long)}:} Inteiros exatos de 32 e 64 bits.
    \item \textbf{\texttt{Double}:} Ponto flutuante de 64 bits.
    \item \textbf{\texttt{Decimal128}:} Decimal de 128 bits com precisão exata para valores monetários e fiscais!
    \item \textbf{\texttt{Date} (\texttt{ISODate}):} Inteiro de 64 bits (ms desde a Unix Epoch).
  \end{itemize}
}

% ------------------------------------------------------------------------------
\begin{frame}[fragile]{Anatomia do \texttt{ObjectId} (12 Bytes) e Sintaxe no Compass (EJSON)}
  \vspace{0.05cm}
  \begin{columns}[t]
    \begin{column}{0.48\textwidth}
      \alert{1. Estrutura dos 12 Bytes do \texttt{ObjectId}}\par
      \vspace{0.08cm}
      \begin{center}
        \tikz[scale=0.68, transform shape]{
          \node[draw=ScoutNavy, fill=ScoutNavy!15, thick, minimum width=2.1cm, minimum height=0.95cm, align=center, font=\scriptsize\bfseries] (b1) at (0,0) {4 Bytes\\Timestamp\\{\tiny\ttfamily 671ec880}};
          \node[draw=ScoutGreen, fill=ScoutGreen!15, thick, minimum width=2.4cm, minimum height=0.95cm, align=center, font=\scriptsize\bfseries, right=0.06cm of b1] (b2) {5 Bytes\\Aleatório (Host)\\{\tiny\ttfamily a1b2c3d4e5}};
          \node[draw=red!80!black, fill=red!10, thick, minimum width=2.0cm, minimum height=0.95cm, align=center, font=\scriptsize\bfseries, right=0.06cm of b2] (b3) {3 Bytes\\Contador\\{\tiny\ttfamily f60001}};
        }
      \end{center}
      \vspace{0.05cm}
      \scriptsize
      \begin{itemize}
        \item Dispensa \texttt{AUTO\_INCREMENT} centralizado: pode ser gerado pelo próprio cliente Python!
        \item Como os 4 primeiros bytes são o \textit{timestamp} Unix, ordenar por \texttt{\_id} já ordena por ordem de criação!
      \end{itemize}
    \end{column}
    \begin{column}{0.48\textwidth}
      \alert{2. Como Inserir no Compass (\textit{Extended JSON})}\par
      \begin{lstlisting}[language=mongosh]
// Na janela Insert Document do Compass, todas
// as chaves exigem aspas duplas ("..."):
{
  "_id": { "$oid": "671ec880a1b2c3d4e5f60001" },
  "sku": "NOTE-ACER-15",
  "nome": "Notebook Gamer Nitro V15",
  "preco": { "$numberDecimal": "4599.90" },
  "estoque": 15,
  "ativo": true
}
      \end{lstlisting}
      \scriptsize
      $\bullet$ \texttt{"\$oid"} $\to$ \texttt{ObjectId}; \texttt{"\$numberDecimal"} $\to$ \texttt{Decimal128}.
    \end{column}
  \end{columns}
\end{frame}

% !TeX root = main.tex
\changecolours{ScoutNavy}
\section[CRUD do Zero: Compass e mongosh]{Aprendendo do Zero: CRUD Passo a Passo no Compass e \texttt{mongosh}}

% ------------------------------------------------------------------------------
\begin{frame}[fragile]{Passo 1: Criando o Banco e Consultando Coleções}
  \vspace{0.05cm}
  \begin{columns}[t]
    \begin{column}{0.48\textwidth}
      \alert{1. Pela Interface Gráfica (MongoDB Compass)}\par
      \itemseps[topsepi=0.2mm,itemsepi=0.5mm,labelsep=1.5mm,leftmargini=2.5mm]
      \scriptsize
      \begin{enumerate}
        \item Conecte em \texttt{mongodb://localhost:27017}.
        \item No menu esquerdo (\textit{Databases}), clique em \textbf{\texttt{+} (Create database)}.
        \item Preencha \textbf{Database Name:} \texttt{mongodb\_books} e \textbf{Collection Name:} \texttt{books}.
        \item \textbf{Consulte a lista de coleções:} Clique sobre \texttt{mongodb\_books} no menu lateral e verifique a coleção \texttt{books} na aba \textbf{Collections}.
      \end{enumerate}
    \end{column}
    \begin{column}{0.48\textwidth}
      \alert{2. Pelo Terminal Interativo (\texttt{mongosh})}\par
      \begin{lstlisting}[language=mongosh]
// Listar todos os bancos existentes:
show dbs

// Selecionar/criar o banco mongodb_books:
use mongodb_books

// Criar explicitamente a colecao books:
db.createCollection("books")

// Consultar a lista de colecoes do banco:
show collections
// Equivalente em JavaScript:
db.getCollectionNames()
      \end{lstlisting}
    \end{column}
  \end{columns}
\end{frame}

% ------------------------------------------------------------------------------
\begin{frame}[fragile]{Passo 2: Inserindo (\textit{Create}) e Consultando (\textit{Read}) o 1.º Documento}
  \vspace{0.05cm}
  \begin{columns}[t]
    \begin{column}{0.48\textwidth}
      \alert{1. Inserção Manual no MongoDB Compass}\par
      \scriptsize
      Em \texttt{books}, clique em \textbf{Add Data $\to$ Insert Document}, cole o JSON e clique em \textbf{Insert}:
      \begin{lstlisting}[language=mongosh]
{
  "isbn": "101",
  "name": "MongoDB 7.0 Guia Pratico",
  "price": 30
}
      \end{lstlisting}
      \scriptsize
      $\bullet$ \textbf{Consultar a lista de documentos:} Na aba \textbf{Documents}, veja que o MongoDB gerou o campo \texttt{\_id} (\texttt{ObjectId}) automaticamente! Alterne entre as visões \textbf{List}, \textbf{JSON} e \textbf{Table}.
    \end{column}
    \begin{column}{0.48\textwidth}
      \alert{2. Equivalente no \texttt{mongosh} (\texttt{insertOne} e \texttt{find})}\par
      \begin{lstlisting}[language=mongosh]
// Inserir 1 documento:
db.books.insertOne({
  isbn: "101",
  name: "MongoDB 7.0 Guia Pratico",
  price: 30
})
// Retorno:
// { acknowledged: true,
//   insertedId: ObjectId("671ec880...") }

// Consultar todos os documentos da colecao:
db.books.find()
      \end{lstlisting}
    \end{column}
  \end{columns}
\end{frame}

% ------------------------------------------------------------------------------
\begin{frame}[fragile]{Passo 3: Subdocumentos Aninhados e Busca com \textit{Dot Notation}}
  \vspace{0.05cm}
  \begin{columns}[t]
    \begin{column}{0.48\textwidth}
      \alert{1. Inserindo Documento com Objeto \texttt{meta}}\par
      \scriptsize
      Em \textbf{Add Data $\to$ Insert Document} no Compass (após inserir, consulte a lista de documentos e expanda \texttt{meta}):
      \begin{lstlisting}[language=mongosh]
{
  "isbn": "201",
  "name": "Modelagem de Documentos NoSQL",
  "meta": {
    "version": "MongoDB 6.0"
  },
  "price": 50
}
      \end{lstlisting}
    \end{column}
    \begin{column}{0.48\textwidth}
      \alert{2. Filtrando no Compass e no \texttt{mongosh}}\par
      \scriptsize
      Na barra \textbf{Filter} do Compass (ou dentro de \texttt{db.books.find(...)}):
      \begin{lstlisting}[language=mongosh]
// 1. Busca exata por atributo de topo:
{ "isbn": "101" }

// 2. Dot Notation em subdocumento:
{ "meta.version": "MongoDB 6.0" }

// 3. Operador de comparacao ($gt = maior que):
{ "price": { "$gt": 35 } }

// 4. Expressao Regular case-insensitive:
{ "name": { "$regex": "mongo", "$options": "i" } }
      \end{lstlisting}
    \end{column}
  \end{columns}
\end{frame}

% ------------------------------------------------------------------------------
\begin{frame}[fragile]{Passo 4: Atualizando (\textit{Update}) e Removendo (\textit{Delete}) Documentos}
  \vspace{0.05cm}
  \begin{columns}[t]
    \begin{column}{0.48\textwidth}
      \alert{1. Atualizando com \texttt{updateOne} e \texttt{\$set}}\par
      \begin{lstlisting}[language=mongosh]
// Atualizar o preco do livro de ISBN "101":
db.books.updateOne(
  { "isbn": "101" },          // Filtro de busca
  { "$set": { "price": 45 } } // Operador atomico
)
// Retorno:
// { acknowledged: true,
//   matchedCount: 1, modifiedCount: 1 }
      \end{lstlisting}
      \scriptsize
      $\bullet$ \textbf{No Compass:} Passe o mouse sobre o documento, clique no lápis (\textbf{Edit Document}), altere \texttt{price} para \texttt{45}, clique em \textbf{Update} e consulte a lista de documentos!
    \end{column}
    \begin{column}{0.48\textwidth}
      \alert{2. Removendo com \texttt{deleteOne} e Lote (\texttt{bulkWrite})}\par
      \begin{lstlisting}[language=mongosh]
// Remover 1 documento especifico:
db.books.deleteOne({ "isbn": "201" })
// Retorno: { acknowledged: true, deletedCount: 1 }

// Operacao em Lote (Batch):
db.books.bulkWrite([
  { insertOne: { document: {
      isbn: "202", name: "MongoDB Guide", price: 40
  }}},
  { updateOne: {
      filter: { isbn: "101" },
      update: { "$set": { price: 49 } }
  }}
])
      \end{lstlisting}
    \end{column}
  \end{columns}
\end{frame}

% !TeX root = main.tex
\changecolours{ScoutGreen}
\section[CRUD com Python (PyMongo)]{Programando com o Driver Oficial Python (\texttt{PyMongo})}

% ------------------------------------------------------------------------------
\begin{frame}[fragile]{Conectando o Python ao MongoDB e Inserindo Documentos}
  \vspace{0.05cm}
  \begin{columns}[t]
    \begin{column}{0.48\textwidth}
      \alert{1. Instalação e Teste de Conexão (\texttt{ping})}\par
      \begin{lstlisting}[language=Python]
# No terminal: pip install pymongo
from pymongo import MongoClient
from pprint import pprint

# Conectar ao servidor local na porta 27017:
uri = "mongodb://localhost:27017/"
client = MongoClient(uri)

# Confirmar conexao com o comando ping:
client.admin.command("ping")
print("Conectado ao MongoDB com sucesso!")

# Selecionar banco 'mongodb_books' e colecao 'books'
books = client.mongodb_books.books
      \end{lstlisting}
    \end{column}
    \begin{column}{0.48\textwidth}
      \alert{2. Inserindo com \texttt{insert\_one} e \texttt{insert\_many}}\par
      \begin{lstlisting}[language=Python]
# Inserir 1 documento (dicionario Python):
book = {
    "isbn": "301",
    "name": "Python and MongoDB",
    "meta": {"version": "MongoDB 7.0"},
    "price": 60
}
res = books.insert_one(book)
print("ID gerado:", res.inserted_id)

# Inserir multiplos documentos em lote:
books.insert_many([
    {"isbn": "302", "name": "Python and MongoDB",
     "meta": {"version": "MongoDB 7.0"}, "price": 60},
    {"isbn": "303", "name": "Advanced MongoDB Techniques",
     "price": 70}
])
      \end{lstlisting}
    \end{column}
  \end{columns}
\end{frame}

% ------------------------------------------------------------------------------
\begin{frame}[fragile]{Consultando Documentos no Python: \texttt{find}, \textit{Dot Notation}, \texttt{\$gt} e \texttt{\$or}}
  \vspace{0.05cm}
  \begin{columns}[t]
    \begin{column}{0.48\textwidth}
      \alert{1. Busca Simples, \textit{Dot Notation} e \texttt{\$regex}}\par
      \begin{lstlisting}[language=Python]
# 1. Buscar pelo nome do livro:
for doc in books.find({"name": "Python and MongoDB"}):
    pprint(doc)

# 2. Buscar dentro do subdocumento 'meta' usando
# Dot Notation e Expressao Regular:
query_regex = {
    "meta.version": {
        "$regex": r".*?g.*?7\.0$",
        "$options": "i"
    }
}
for doc in books.find(query_regex):
    pprint(doc)
      \end{lstlisting}
    \end{column}
    \begin{column}{0.48\textwidth}
      \alert{2. Comparação (\texttt{\$gt}) e Operadores Lógicos}\par
      \begin{lstlisting}[language=Python]
# 3. Livros com preco maior que 50 ($gt):
for doc in books.find({"price": {"$gt": 50}}):
    pprint(doc)

# 4. Condicao logica AND (multiplas chaves):
q_and = {"name": "Advanced MongoDB Techniques",
         "price": 70}
pprint(books.find_one(q_and))

# 5. Condicao logica OR ($or):
q_or = {
    "$or": [
        {"name": "Advanced MongoDB Techniques"},
        {"isbn": "301"}
    ]
}
for doc in books.find(q_or):
    pprint(doc)
      \end{lstlisting}
    \end{column}
  \end{columns}
\end{frame}

% ------------------------------------------------------------------------------
\begin{frame}[fragile]{Atualizando (\texttt{update\_one}) e Removendo (\texttt{delete\_one}) no Python}
  \vspace{0.05cm}
  \begin{columns}[t]
    \begin{column}{0.48\textwidth}
      \alert{1. Atualizando com \texttt{update\_one} e \texttt{\$set}}\par
      \begin{lstlisting}[language=Python]
# Atualizar o preco do livro "Advanced MongoDB
# Techniques" para 75:
update_result = books.update_one(
    {"name": "Advanced MongoDB Techniques"},
    {"$set": {"price": 75}}
)

print("Encontrados:", update_result.matched_count)
print("Modificados:", update_result.modified_count)

# Buscar o documento atualizado para confirmar:
updated_doc = books.find_one(
    {"name": "Advanced MongoDB Techniques"}
)
pprint(updated_doc)
      \end{lstlisting}
    \end{column}
    \begin{column}{0.48\textwidth}
      \alert{2. Removendo com \texttt{delete\_one}}\par
      \begin{lstlisting}[language=Python]
# Remover o livro de ISBN '303':
isbn_to_delete = "303"
delete_result = books.delete_one(
    {"isbn": isbn_to_delete}
)
print("Removidos:", delete_result.deleted_count)

# Listando colecoes do banco via Python:
db = client.mongodb_books
print("Colecoes:", db.list_collection_names())
      \end{lstlisting}
      \scriptsize
      $\bullet$ \textbf{Sincronia com o Compass:} Clique em \textbf{Refresh} na aba \textbf{Documents} do Compass para ver as alterações!
    \end{column}
  \end{columns}
\end{frame}

% !TeX root = main.tex
\changecolours{ScoutNavy}
\section[Modelagem: Embedding vs. Referencing]{Teoria de Schema Design: \textit{Embedding}, \textit{Referencing} e Padrões}

% ------------------------------------------------------------------------------
\twocolframe[title={Princípios de Schema Design no MongoDB},leftcol=7.0cm,rightcol=7.0cm]{%
  \alert{Três Estratégias Fundamentais}\par
  \itemseps[topsepi=0.2mm,itemsepi=0.5mm,labelsep=1.5mm,leftmargini=2.5mm]
  \scriptsize
  \begin{enumerate}
    \item \textbf{\textit{Embedded Data} (Dados Incorporados):} Aninhar objetos relacionados dentro do mesmo documento (ex.: um livro com o subdocumento \texttt{author}, ou uma cliente com a lista de \texttt{enderecos}).
    \item \textbf{\textit{References} (Referências):} Guardar o \texttt{ObjectId} de um documento em outra coleção (ex.: coleção \texttt{courses} apontando para \texttt{professor: ObjectId(...)}).
    \item \textbf{\textit{Denormalization} (Desnormalização):} Copiar campos específicos (como \texttt{authorName}) para evitar \texttt{\$lookup} em leituras críticas.
  \end{enumerate}
}{%
  \alert{A Métrica-Chave: \textit{Read-Write Ratio}}\par
  \itemseps[topsepi=0.2mm,itemsepi=0.5mm,labelsep=1.5mm,leftmargini=2.5mm]
  \scriptsize
  \begin{itemize}
    \item \textbf{Proporção Leitura vs. Escrita (\textit{Read-Write Ratio}):}
      \begin{itemize}\scriptsize
        \item \textbf{Cargas Intensivas em Leitura (\textit{Read-Heavy}):} Catálogos de produtos e vitrines têm muito mais leituras que escritas $\to$ favorecem \textbf{\textit{Embedding}} e \textbf{Desnormalização} (1 única busca traz tudo).
        \item \textbf{Cargas Intensivas em Escrita (\textit{Write-Heavy}):} Eventos frequentes ou entidades compartilhadas que mudam o tempo todo $\to$ favorecem \textbf{\textit{Referencing}}.
      \end{itemize}
    \item \textbf{Regra de Ouro:} \textit{``Dados lidos juntos pela aplicação devem ser armazenados juntos!''}
  \end{itemize}
}

% ------------------------------------------------------------------------------
\begin{frame}{Visualizando as Duas Estratégias de Modelagem}
  \vspace{0.15cm}
  \makebox[\textwidth][c]{%
    \begin{tikzpicture}
      % Bloco 1: Embedding (Esquerda)
      \node[draw=ScoutGreen, fill=ScoutGreen!10, very thick, rounded corners=6pt, minimum width=6.8cm, minimum height=5.0cm] (emb) at (0,0) {};
      \node[font=\bfseries\sffamily\footnotesize, text=ScoutGreen] at (0, 2.1) {1. Incorporação (\textit{Embedding} / Desnormalizado)};
      \node[draw=ScoutGreen!80!black, fill=white, thick, rounded corners=4pt, text width=6.1cm, inner sep=6pt, align=left, font=\ttfamily\fontsize{6.2pt}{7.8pt}\selectfont] at (0, -0.25) {%
        \textbf{\sffamily\color{ScoutGreen}Coleção: clientes (1 Único Documento)}\\[2pt]
        \{\\
        \hspace*{2mm}"\_id": \{ "\$oid": "671ec880...0010" \},\\
        \hspace*{2mm}"nome": "Ana Costa",\\
        \hspace*{2mm}"cpf": "123.456.789-00",\\
        \hspace*{2mm}\textbf{"enderecos": [} \textcolor{ScoutGreen}{// Subdocumentos!}\\
        \hspace*{4mm}\{ "apelido": "Casa", "rua": "Av. Chermont" \},\\
        \hspace*{4mm}\{ "apelido": "IFCE", "rua": "CE-187, Taua" \}\\
        \hspace*{2mm}\textbf{]}\\
        \}%
      };

      % Bloco 2: Referencing (Direita)
      \node[draw=ScoutNavy, fill=ScoutNavy!10, very thick, rounded corners=6pt, minimum width=6.8cm, minimum height=5.0cm] (ref) at (7.3,0) {};
      \node[font=\bfseries\sffamily\footnotesize, text=ScoutNavy] at (7.3, 2.1) {2. Referenciamento (\textit{Referencing} / Normalizado)};

      \node[draw=ScoutNavy!80!black, fill=white, thick, rounded corners=4pt, text width=6.1cm, inner sep=5.5pt, align=left, font=\ttfamily\fontsize{6.2pt}{7.6pt}\selectfont] (docA) at (7.3, 0.85) {%
        \textbf{\sffamily\color{ScoutNavy}Coleção: avaliacoes}\\[2pt]
        \{\\
        \hspace*{2mm}"\_id": \{ "\$oid": "671ec990...0999" \},\\
        \hspace*{2mm}\textbf{\color{red!80!black}"produto\_id": \{ "\$oid": "671ec880...0001" \}},\\
        \hspace*{2mm}"autor": "Lucas", "nota": 5\\
        \}%
      };

      \node[draw=ScoutNavy!80!black, fill=white, thick, rounded corners=4pt, text width=6.1cm, inner sep=5.5pt, align=left, font=\ttfamily\fontsize{6.2pt}{7.6pt}\selectfont] (docB) at (7.3, -1.55) {%
        \textbf{\sffamily\color{ScoutNavy}Coleção: produtos}\\[2pt]
        \{\\
        \hspace*{2mm}\textbf{\color{red!80!black}"\_id": \{ "\$oid": "671ec880...0001" \}},\\
        \hspace*{2mm}"sku": "NOTE-ACER-15", "preco": 4599.90\\
        \}%
      };

      \draw[->, ultra thick, red!80!black] ([xshift=-1.4cm]docA.south) -- node[right=4pt, font=\bfseries\sffamily\scriptsize, text=red!80!black] {Referência por \texttt{ObjectId}} ([xshift=-1.4cm]docB.north);
    \end{tikzpicture}%
  }
\end{frame}

% ------------------------------------------------------------------------------
\begin{frame}{O Anti-Padrão \textit{Unbounded Arrays} e os Padrões de Projeto}
  \vspace{0.05cm}
  \begin{columns}[t]
    \begin{column}{0.46\textwidth}
      \alert{Por que Evitar Arrays Ilimitados?}\par
      \itemseps[topsepi=0.2mm,itemsepi=0.4mm,labelsep=1.5mm,leftmargini=2.5mm]
      \scriptsize
      \begin{itemize}
        \item Embutir listas sem fim (ex.: 50.000 avaliações no produto) estoura o limite de \textbf{16 MB} (\texttt{BSONObjectTooLarge}) e degrada o \textit{WiredTiger}.
        \item \textbf{Cardinalidade:} \textbf{1:1 / 1:Few ($<50$)} $\to$ \textit{Embedding}; \textbf{1:Many} $\to$ \textit{Referencing} / \textit{Subset}; \textbf{1:Milhões} $\to$ Ref. no filho.
      \end{itemize}
    \end{column}
    \begin{column}{0.50\textwidth}
      \alert{Padrões de Modelagem (\textit{Design Patterns})}\par
      \vspace{0.03cm}
      \scriptsize
      \renewcommand{\arraystretch}{1.08}
      \begin{tabular}{p{1.65cm} p{2.65cm}}
        \toprule
        \textbf{Padrão} & \textbf{Aplicação Prática} \\
        \midrule
        \textbf{Polymorphic} & Categorias distintas na mesma coleção. \\
        \textbf{Subset} & Últimas avaliações no produto. \\
        \textbf{Extended Ref} & Campos do cliente/item no pedido. \\
        \textbf{Computed} & Média e total pré-calculados. \\
        \bottomrule
      \end{tabular}
    \end{column}
  \end{columns}
\end{frame}

% !TeX root = main.tex
\changecolours{ScoutGreen}
\section[Prática no ShopGrids: Compass e Python]{Modelando o E-Commerce ShopGrids no Compass e no Python}

% ------------------------------------------------------------------------------
\begin{frame}[fragile]{Coleção \texttt{produtos}: \textit{Polymorphic Pattern} e Validação}
  \vspace{0.05cm}
  \begin{columns}[t]
    \begin{column}{0.48\textwidth}
      \alert{1. Documento 1: Eletrônicos (Compass)}\par
      \begin{lstlisting}[language=mongosh]
{
  "_id": { "$oid": "671ec880a1b2c3d4e5f60001" },
  "sku": "NOTE-ACER-15",
  "nome": "Notebook Gamer Nitro V15",
  "categoria": "Eletronicos",
  "preco": { "$numberDecimal": "4599.90" },
  "estoque": 15,
  "especificacoes": {
    "cpu": "Intel Core i7-13620H",
    "ram_gb": 16,
    "gpu": "RTX 4050 6GB"
  },
  "garantia_meses": 12
}
      \end{lstlisting}
    \end{column}
    \begin{column}{0.48\textwidth}
      \alert{2. Documento 2: Vestuário + Validação}\par
      \begin{lstlisting}[language=mongosh]
{
  "_id": { "$oid": "671ec880a1b2c3d4e5f60002" },
  "sku": "CAM-TECH-ADS",
  "nome": "Camiseta Dev NoSQL IFCE",
  "categoria": "Vestuario",
  "preco": { "$numberDecimal": "69.90" },
  "estoque": 80,
  "especificacoes": {
    "tecido": "100% Algodao Egipcio",
    "cor": "Preta",
    "tamanhos": ["P", "M", "G", "GG"]
  }
}
      \end{lstlisting}
      \scriptsize
      $\bullet$ \textbf{Governança (\texttt{\$jsonSchema} / Aba Validation):} Exige \texttt{sku}, \texttt{nome} e \texttt{preco > 0}, deixando \texttt{especificacoes} livre!
    \end{column}
  \end{columns}
\end{frame}

% ------------------------------------------------------------------------------
\begin{frame}[fragile]{Coleções \texttt{avaliacoes} e \texttt{pedidos}: \textit{Subset}, \textit{Computed} e \textit{Extended Reference}}
  \vspace{0.02cm}
  \begin{columns}[t]
    \begin{column}{0.48\textwidth}
      \alert{1. \textit{Subset} + \textit{Computed} (\texttt{avaliacoes})}\par
      \begin{lstlisting}[language=mongosh]
// Em produtos (Vitrine em 1 Query!):
{
  "_id": { "$oid": "671ec880a1b2c3d4e5f60001" },
  "sku": "NOTE-ACER-15",
  "resumo_avaliacoes": { "media": 5.0, "total": 2 },
  "ultimas_avaliacoes": [
    { "autor": "Lucas", "nota": 5, "txt": "Otimo!" }
  ]
}
// Em avaliacoes (Historico Referenciado):
{
  "_id": { "$oid": "671ec990a1b2c3d4e5f60999" },
  "produto_id": { "$oid": "671ec880a1b2c3d4e5f60001" },
  "autor": "Lucas", "nota": 5, "txt": "Otimo!"
}
      \end{lstlisting}
    \end{column}
    \begin{column}{0.48\textwidth}
      \alert{2. \textit{Extended Ref.} + \textit{Snapshot} (\texttt{pedidos})}\par
      \begin{lstlisting}[language=mongosh]
{
  "_id": { "$oid": "671eca00a1b2c3d4e5f60777" },
  "codigo_pedido": "PED-2026-0001",
  "cliente": {
    "cliente_id": { "$oid": "671ec880...0010" },
    "nome": "Ana Costa", "cpf": "123.456.789-00"
  },
  "itens": [{
    "produto_id": { "$oid": "671ec880...0001" },
    "sku": "NOTE-ACER-15",
    "preco_unit": { "$numberDecimal": "4599.90" },
    "quantidade": 1
  }],
  "valor_total": { "$numberDecimal": "4599.90" },
  "status": "PAGO"
}
      \end{lstlisting}
    \end{column}
  \end{columns}
\end{frame}

% ------------------------------------------------------------------------------
\begin{frame}{Arquitetura Completa das 4 Coleções do E-Commerce ShopGrids}
  \vspace{0.15cm}
  \makebox[\textwidth][c]{%
    \begin{tikzpicture}
      % Colecao Produtos
      \node[draw=ScoutGreen, fill=ScoutGreen!12, very thick, rounded corners=5pt, text width=5.4cm, minimum height=2.1cm, inner sep=6pt, align=left, font=\sffamily\scriptsize] (prod) at (0, 1.45) {%
        \textbf{\color{ScoutGreen}\footnotesize Coleção: \texttt{produtos}}\\[2pt]
        $\bullet$ \texttt{\_id: ObjectId(...)}\\
        $\bullet$ \texttt{sku, nome, preco (Decimal128)}\\
        $\bullet$ \textbf{Polymorphic 1:1:} \texttt{especificacoes: \{...\}}\\
        $\bullet$ \textbf{Subset + Computed:} \texttt{ultimas\_avaliacoes}%
      };

      % Colecao Avaliacoes
      \node[draw=ScoutNavy, fill=ScoutNavy!12, very thick, rounded corners=5pt, text width=5.4cm, minimum height=1.7cm, inner sep=6pt, align=left, font=\sffamily\scriptsize] (aval) at (0, -1.45) {%
        \textbf{\color{ScoutNavy}\footnotesize Coleção: \texttt{avaliacoes}}\\[2pt]
        $\bullet$ \texttt{\_id: ObjectId(...)}\\
        $\bullet$ \textbf{Ref 1:Many:} \texttt{produto\_id: ObjectId(...)}\\
        $\bullet$ \texttt{autor, nota, txt}%
      };

      % Colecao Clientes
      \node[draw=ScoutGreen, fill=ScoutGreen!12, very thick, rounded corners=5pt, text width=5.6cm, minimum height=2.1cm, inner sep=6pt, align=left, font=\sffamily\scriptsize] (cli) at (8.0, 1.45) {%
        \textbf{\color{ScoutGreen}\footnotesize Coleção: \texttt{clientes}}\\[2pt]
        $\bullet$ \texttt{\_id: ObjectId(...)}\\
        $\bullet$ \texttt{nome, email, cpf}\\
        $\bullet$ \textbf{Embedded 1:Few:}\\
        \hspace*{3mm}\texttt{enderecos: [ \{Casa\}, \{IFCE\} ]}%
      };

      % Colecao Pedidos
      \node[draw=red!80!black, fill=red!10, very thick, rounded corners=5pt, text width=5.6cm, minimum height=1.7cm, inner sep=6pt, align=left, font=\sffamily\scriptsize] (ped) at (8.0, -1.45) {%
        \textbf{\color{red!80!black}\footnotesize Coleção: \texttt{pedidos}}\\[2pt]
        $\bullet$ \texttt{\_id: ObjectId(...), status}\\
        $\bullet$ \textbf{Extended Ref:} \texttt{cliente: \{cliente\_id, nome\}}\\
        $\bullet$ \textbf{Snapshot:} \texttt{itens: [ \{produto\_id, preco\_unit\} ]}%
      };

      \draw[->, ultra thick, ScoutNavy] (aval.north) -- node[right=2pt, font=\scriptsize\bfseries\sffamily] {\texttt{produto\_id}} (prod.south);
      \draw[->, ultra thick, red!80!black] (ped.north) -- node[right=2pt, font=\scriptsize\bfseries\sffamily] {\texttt{cliente\_id}} (cli.south);
      \draw[->, ultra thick, red!80!black] (ped.west) -- node[above, sloped, font=\scriptsize\bfseries\sffamily] {\texttt{itens.produto\_id}} (prod.east);
    \end{tikzpicture}%
  }
\end{frame}

% ------------------------------------------------------------------------------
\begin{frame}{Checklist de Decisão e Roteiro Passo a Passo (Compass + Python)}
  \vspace{0.1cm}
  \makebox[\textwidth][c]{%
    \scriptsize
    \renewcommand{\arraystretch}{1.35}
    \begin{tabular}{p{5.4cm} p{3.2cm} p{3.2cm}}
      \toprule
      \textbf{Pergunta de Projeto} & \textbf{\textit{Embedding}} & \textbf{\textit{Referencing}} \\
      \midrule
      1. Dados lidos juntos na mesma tela? & \textbf{Sim} (1 única query) & Não (telas separadas) \\
      2. Qual a cardinalidade máxima? & \textbf{1:1} ou \textbf{1:Few} ($<50$) & \textbf{1:Many} / \textbf{1:Milhões} \\
      3. Array pode crescer sem limite? & Não (tamanho fixo) & \textbf{Sim} (risco de 16 MB) \\
      4. Dado filho consultado isoladamente? & Raramente & \textbf{Sim, sempre} \\
      5. Qual o \textit{Read-Write Ratio}? & \textbf{Alta Leitura} (\textit{Read-Heavy}) & \textbf{Alta Escrita} (\textit{Write-Heavy}) \\
      6. Congelar valor fiscal histórico? & \textbf{Sim (\textit{Snapshot})} & Não (aponta para atual) \\
      \bottomrule
    \end{tabular}%
  }
  \vspace{0.15cm}
  \centering{\scriptsize \textbf{Sequência da Prática:} \textbf{Etapa 1:} CRUD do Zero em \texttt{mongodb\_books} (Compass + \texttt{mongosh}) $\to$ \textbf{Etapa 2:} CRUD com Python (\texttt{PyMongo}) $\to$ \textbf{Etapa 3:} Modelagem manual do 1.º registro de cada coleção em \texttt{ecommerce\_shopgrids} no Compass $\to$ \textbf{Etapa 4:} Povoamento em lote via Compass/Python.}
\end{frame}


% ------------------------------------------------------------------------------
% Encerramento
% ------------------------------------------------------------------------------
\changecolours[inverse]{ScoutGreen}
\begin{frame}[plain]
  \vspace{2cm}
  \centering
  {\Huge \textbf{Dúvidas \& Prática em Laboratório}}\par
  \vspace{0.8cm}
  {\large IFCE Campus Tauá -- ADS16 (Modelo Orientado a Documentos)}\par
  \vspace{0.3cm}
  {\normalsize Do CRUD Básico ao Python (\texttt{PyMongo}) e Modelagem BSON no MongoDB}
\end{frame}

\end{document}
